\subsection{Design by Contract and Specification-Driven Development}
\label{subsec:rw2}

The notion that a program component should be governed by a precise, machine-checkable agreement between supplier and client predates the present interest in agentic software construction by several decades, and it is from this lineage that the formal-contract layer proposed in this article descends. The semantic foundations were laid by Hoare's axiomatic treatment of program correctness, in which a triple relates a precondition, a command, and a postcondition~\cite{hoare1969axiomatic}, and by Dijkstra's predicate-transformer calculus, which renders the weakest precondition of a statement with respect to a desired postcondition computable~\cite{dijkstra1976discipline}. These two results are the technical bedrock on which both deductive verification and the compositional change-propagation analysis used later in this article rest: the weakest-precondition transformer is precisely the mechanism by which a change to a callee's contract can be pushed back through its call sites.

Meyer reified these ideas into an engineering discipline he named Design by Contract (DbC), embedding preconditions, postconditions, and class invariants as first-class, executable constructs in the Eiffel language~\cite{meyer1992dbc,meyer1997oosc}. The decisive move was to treat the contract not as external documentation but as a durable program artefact with defined obligations on each party and defined blame when an obligation is violated. This framing---contract as the authoritative statement of intent, code as its fulfillment---directly anticipates the economic inversion this article builds upon, in which a cheaply regenerable implementation becomes a build output and the contract becomes the scarce, durable asset. DbC's principal limitation, however, is that its assertions are checked only at runtime against the inputs that happen to occur; the guarantee is sampled, not universal.

The behavioural interface specification tradition closes part of this gap by enriching contracts to a point at which static proof becomes tractable. The Java Modeling Language (JML) marries Eiffel-style contracts with the model-based specification heritage of Larch and elements of the refinement calculus, yielding a notation expressive enough for full functional specification of Java classes~\cite{leavens2006jml}. OpenJML discharges such specifications by Extended Static Checking (ESC), translating a method together with its contract into verification conditions checked by an SMT solver, thereby establishing correctness over \emph{all} inputs rather than a sampled subset~\cite{cok2011openjml,cok2014openjml}. Spec\#~\cite{barnett2005specsharp} pursued the same agenda for C\#, demonstrating that a contract layer can be retrofitted onto a mainstream object-oriented language with a verifying compiler. Underpinning the static-proof tradition is the refinement calculus of Back and Morgan~\cite{back1998refinement,morgan1994programming}, which formalizes the stepwise transformation of an abstract specification into executable code with each step preserving correctness---the correctness-by-construction ideal that the synthesis-and-verification loop in this article approximates pragmatically, replacing manual refinement steps with agent-generated candidates that are filtered by a verifier.

A persistent obstacle to all of the above is the cost of authoring formal contracts in the first place, and a substantial body of work addresses this by \emph{inferring} specifications rather than demanding them. Dynamic detection, exemplified by Daikon, observes executions and reports likely invariants~\cite{ernst2007daikon}, while Houdini infers candidate annotations and retains those an underlying checker can prove~\cite{flanagan2001houdini}. More recently, large language models have been applied to specification synthesis: SpecGen drives an LLM through conversational refinement and mutation, selecting candidates a verifier accepts~\cite{ma2025specgen}, and related efforts translate natural-language intent into temporal or behavioural specifications~\cite{cosler2023nl2spec}. The author's own prior work fits squarely in this inference tradition, recovering JML by AST heuristics and showing that such specifications materially improve downstream artefacts and admit compositional refinement across method boundaries~\cite{edmonds_article1,edmonds_article3}.

What this collected literature does not supply is the missing element for the emerging spec-driven, agentic lifecycle. The classical DbC and JML traditions assume a human author who can and will write formal contracts, an assumption the natural-language–first agentic workflow explicitly contradicts; inference work, conversely, recovers contracts that faithfully describe what code \emph{does} rather than what it \emph{should}, and---crucially---does not address \emph{who} owns the resulting contract when that same contract must serve as the oracle against which generated code is judged. None of these strands treats the specification simultaneously as the input the agent builds against and the oracle it is verified against, nor do they confront the brownfield reality in which a contract must be triangulated from code, issue tickets, and documentation with explicit provenance, conflicts surfaced as candidate bugs rather than suppressed. This article positions its machine-checkable contract layer, its independent-authority and ratification constraints, and its compositional change-propagation analysis precisely at that gap, promoting the inference and verification machinery surveyed here into a governed lifecycle pivot rather than an isolated tool.
